\documentclass[10pt,a4paper]{article}
\usepackage[margin=2cm]{geometry}
\usepackage{amsmath,amssymb,booktabs,enumitem,hyperref,parskip,xcolor}
\setlist{nosep,leftmargin=1.4em}
\hypersetup{colorlinks=true,urlcolor=blue!60!black,linkcolor=black}
\newcommand{\CE}{\mathrm{CE}}
\newcommand{\fw}{f_w}
\newcommand{\fhat}{\hat f_t}
\newcommand{\ftil}{\tilde f_\tau}
\newcommand{\pbar}{\bar p}
\newcommand{\Q}[1]{\par\medskip\noindent\textbf{Q: #1}\par\noindent}
\title{\vspace{-1.5cm}Group 3 proposal: theory, speaker notes and Q\&A\\[0.2em]
\large Uncertainty-Aware Weak-to-Strong Generalization (deck: \texttt{G3\_proposal\_slides.pdf})}
\author{\normalsize Trustworthy AI, NYCU \quad$\cdot$\quad proposal talk Wed 18 Nov 2026, A304: 10 min + 5 min Q\&A}
\date{}
\begin{document}
\maketitle
\vspace{-1.2cm}

\section{The proposal in five sentences}
A02 (Burns et al., 2024) trains a strong pretrained student on the labels of a weak teacher and adds a
\emph{confidence loss} that makes the student also trust its own hardened prediction $\fhat$.
The paper judges the result by accuracy only.
U09 (Thuy \& Benoit, 2024) shows how to estimate uncertainty (MC Dropout, ensembles) and use it to
\emph{reject} doubtful predictions, but only after training and with clean labels.
We put U09's rejection step \emph{inside} A02's loss: the student hardens its own prediction only where its
model uncertainty is low.
We then ask whether this student is better at telling a human which predictions to check, at equal accuracy.

\section{Theory: how and why to modify $\fhat$}

\subsection{What the A02 loss does}
\[
L_{\text{conf}}(f)=(1-\alpha)\,\CE(f(x),\fw(x))+\alpha\,\CE(f(x),\fhat(x))
=\CE\big(f(x),\;(1-\alpha)\fw(x)+\alpha\fhat(x)\big)
\]
(App.\ A.4, Eqs.\ 1--2; the two forms are equal because cross-entropy is linear in its target).
$f(x)$ is the student's predicted distribution, $\fw(x)$ the teacher's, and
$\fhat(x)=I[f(x)>t]$ the student's own prediction turned into a hard 0/1 label.
$\alpha$ is a fixed weight (0.5, or 0.75 for the largest students) warmed up linearly from 0 over the first
20\,\% of training. In the paper's binary tasks $t$ is set per batch so that exactly half of the batch is
labelled positive (a class-prior constraint). With 10 classes the natural version is
$\fhat=\text{one-hot}(\arg\max f(x))$; the paper's footnote says the prior should then be computed from the
weak predictions.

\textbf{Gradient view (the key to everything below).} For a fixed target $q$, the gradient of
$\CE(f,q)$ with respect to the logits is $f-q$. So the confidence loss pushes the logits along
\[
(1-\alpha)\,\underbrace{(f-\fw)}_{\text{imitate the teacher}}+\;\alpha\,\underbrace{(f-\fhat)}_{\text{sharpen your own answer}} .
\]
The second term moves probability mass towards the class the student already prefers: it is entropy
minimisation / self-training (the paper cites Grandvalet \& Bengio 2004, Lee 2013, Reed et al.\ 2014).
It helps when the student's own answer is better than the teacher's label, which is the whole point of
weak-to-strong. It hurts when the student has already \emph{copied} a teacher mistake: then
$\fhat=\fw$=the wrong class, and both terms push confidence in the wrong class up. In Tom\'as's earlier pilot
(AlexNet $\to$ DINO ViT-B/8) copied errors were indeed more confident than the student's other errors
(0.71 vs 0.61 with the confidence loss).

\subsection{Why changing $\alpha$ per image is not really a new idea (the professor's point)}
Our first idea (method C) was $\alpha_i=\alpha(1-u_i)$ with $u_i\in[0,1]$ the uncertainty of image $i$.
Substituting into the target:
\[
(1-\alpha_i)\fw+\alpha_i\fhat=(1-\alpha)\,\fw+\alpha\,\big[(1-u_i)\,\fhat+u_i\,\fw\big].
\]
So ``per-image $\alpha$'' \emph{is} a modification of $\fhat$: it replaces $\fhat$ by a blend that
\textbf{falls back to the teacher's label when the student is uncertain}. That answers the
``same loss?'' question in our notes: yes, it is the same loss family, and seen this way the choice is not
obviously good, because uncertain images are plausibly the ones where the teacher is wrong more often
(we will measure this: hypothesis H3). The cleaner question is the professor's: \emph{what should the
self-target be?}

\subsection{Our self-target}
For each training image, run $T$ stochastic forward passes of the student head (MC Dropout, e.g.\ $T=10$,
no gradient): $p_1,\dots,p_T$, mean $\pbar=\frac1T\sum_tp_t$. U09's decomposition (also UQ deck, MC Dropout
section):
\[
\underbrace{H(\pbar)}_{\text{total}}=\underbrace{\tfrac1T\textstyle\sum_tH(p_t)}_{\text{data (aleatoric)}}+
\underbrace{H(\pbar)-\tfrac1T\textstyle\sum_tH(p_t)}_{\text{model (epistemic)}=:u(x)} .
\]
\[
\ftil(x)=\begin{cases}\text{one-hot}(\arg\max\pbar(x)) & u(x)\le\tau \quad\text{(certain: harden, as A02)}\\
\pbar(x) & u(x)>\tau\quad\text{(uncertain: abstain)}\end{cases}
\quad L=\CE\big(f,(1-\alpha)\fw+\alpha\ftil\big).
\]
$\tau$ is the $q$-quantile of $u$ in the batch, so a fraction $q$ of each batch is hardened. $q=1$ gives A02
(with the MC mean instead of a single pass); $q=0$ never hardens. $\alpha$ and its warm-up stay as in A02.

\textbf{Why this is a sensible reading of the professor's hint.}
\begin{itemize}
\item In A02 the only place where the student's \emph{own belief} enters the loss is $\fhat$, and it enters as
a point estimate with all uncertainty thrown away. Making $\fhat$ uncertainty-aware is the minimal change.
\item A02 already has a threshold $t$ inside $\fhat$, set per batch. We keep that structure and change what
is thresholded: the model uncertainty instead of the raw score.
\item By the gradient view, on abstained images the self-term is $\alpha(f-\pbar)\approx0$: nothing is
sharpened, the student stays as unsure as it is. On hardened images it behaves exactly like A02. So
confidence is amplified only where the student's belief is stable across dropout masks.
\item It is U09's ``classification with rejection'', applied to the student's own pseudo-labels at training
time instead of to predictions at test time. $q$ plays the role of coverage.
\end{itemize}

\textbf{Why model uncertainty and not total entropy.} Total entropy is high on ambiguous images (data
uncertainty) too; hardening there is harmless. Model uncertainty is high where heads that fit the weak
labels equally well \emph{disagree}, which is where the weak labels fight the pretrained features. That is
our candidate signal for ``the teacher may be wrong here''. It is a hypothesis, so the ablation compares
total / data / model and a \emph{random} gate with the same $q$.

\textbf{Variants (ablations D$'$).} (i) soft blend $\ftil=(1-\tilde u)\,\text{one-hot}+\tilde u\,\pbar$ with
$\tilde u$ the batch rank of $u$, instead of a hard gate; (ii) ensemble of 5 heads instead of MC Dropout;
(iii) method C, i.e.\ fallback to the teacher.

\subsection{What would be an interesting result}
\begin{enumerate}
\item \textbf{D beats B and the random gate on error detection / accuracy--coverage at equal accuracy.}
Positive result: uncertainty-aware self-training gives a more trustworthy weak-to-strong student.
\item \textbf{D $\approx$ random gate.} Then model uncertainty carries no information about teacher errors in this
setting (H3 fails): a clear negative finding about MC Dropout on frozen features, and the gain of C in the
pilot was only due to less hardening on average.
\item \textbf{Accuracy falls as $q$ decreases.} Then hardening everywhere is what produces the PGR gain, and there is
a measurable trade-off between accuracy and uncertainty quality: plot both against $q$.
\end{enumerate}
In all three cases the figure ``accuracy and error-detection AUROC as a function of $q$'' is the centre of the poster.

\subsection{Link to the course}
\begin{itemize}
\item \textbf{Intro-TAI requirements table:} we target \emph{uncertainty awareness}, \emph{reliability} and
\emph{human oversight}; accuracy is a constraint, not the goal. Lifecycle view: we intervene at the
\emph{model} (training) stage to improve a \emph{decision} stage policy (defer or predict).
\item \textbf{UQ deck, what uncertainty enables (pp.\ 16--25):} safer prediction,
$\text{defer}(x)=1\{U(x)>\tau\}$. Our gate is the same rule used at learning time.
\item \textbf{UQ deck, sources (aleatoric vs epistemic) and methods (p.\ 35, tables pp.\ 37, 98):} MC Dropout and
ensembles are the ``generate multiple predictions'' route; their disagreement is a \emph{proxy} for epistemic
uncertainty, not a guarantee.
\item \textbf{UQ deck, evaluating uncertainty (pp.\ 99--105):} calibration (ECE, reliability diagram), ``does
uncertainty rank risk?'' (selective prediction, risk--coverage curve, AURC), ``does it survive shift?''
(named corruptions: the optional CIFAR-10-C test).
\item \textbf{UQ deck, from scores to decisions (pp.\ 107--113):} a validated score is not yet a policy; deferring
the most uncertain predictions to a human is the decision our evaluation targets.
\item \textbf{U09 / XAI:} U09 presents uncertainty as a local, model-specific explanation (``why should I not trust
this prediction''), which is how we link to the XAI part of the course.
\end{itemize}

\section{Speaker notes (about 10 minutes; suggested split, change freely)}
\textbf{Slide 1, title (20 s).} We are Group 3. Our project combines two of our midterm papers: weak-to-strong
generalization and explainability through uncertainty. In one line: we let a strong student abstain from its
own pseudo-labels during training.

\textbf{Slide 2, issue and motivation (70 s).} The trustworthy AI issue is reliability and human oversight when
the training labels come from a weaker supervisor. A strong pretrained model learns from a weaker model's
labels, because no ground truth is available. It becomes more accurate than its teacher, but it also copies
some of the teacher's mistakes, and it is confident about them, so nobody is asked to check them. What is
needed is a model that sends the right predictions to a human. Accuracy alone does not give that.

\textbf{Slide 3, A02 (80 s).} Our first key paper studies exactly this setting. A strong student is trained on a
weak supervisor's labels. Their best method is the auxiliary confidence loss: the training target is a mix of
the teacher's label and the student's own hardened prediction, f-hat. Alpha controls the mix. This lets the
student disagree with its teacher and it increases the performance gap recovered. Two things remain open. The
paper only reports accuracy. And f-hat is a point estimate that is hardened on every example. If the student
has already copied a teacher error, f-hat equals the wrong label and the loss reinforces it. The official code
is public, including the vision experiment, so the baseline is reproducible.

\textbf{Slide 4, U09 (70 s).} Our second key paper gives the missing tool. Stage one: estimate uncertainty with
several stochastic forward passes, MC Dropout or deep ensembles, and split total uncertainty into a data part
and a model part. Stage two: classification with rejection: the most uncertain cases go to a human, and
accuracy on the rest goes up, even under distribution shift. Its limits for us: it assumes clean labels, and
uncertainty is only used after training. We re-implement it; MC Dropout is a few lines.

\textbf{Slide 5, research question and contribution (60 s).} Our question: if we make the self-training target
uncertainty-aware, does the student become more reliable about which predictions to trust, without losing
accuracy? On the methodological side we propose an uncertainty-gated self-target, of which A02 is a special
case, and we evaluate weak-to-strong learning by reliability, not only accuracy. On the practical side this is
a training recipe for any pipeline that replaces a weak labeller with a pretrained model and has a limited
human review budget. Note that we do not change alpha; we change where the student is allowed to trust itself.

\textbf{Slide 6, method (100 s).} Three steps. One: for every training image we run T dropout passes of the
student and compute the mean prediction and the model uncertainty, which is the disagreement between passes.
Two: we harden the prediction only if the model uncertainty is below a threshold tau; otherwise we keep the
soft mean prediction. Three: the A02 loss is unchanged, with our target in place of f-hat. Tau is a batch
quantile, so a fraction q of each batch is hardened; q equal to one is A02. Why should it help? The gradient
of cross-entropy is prediction minus target. A hard target pushes confidence up. A soft target equal to the
student's own mean pushes almost nowhere. So confidence only grows where the student's belief is stable. This
is U09's rejection, applied to the student's own pseudo-labels.

\textbf{Slide 7, data and models (70 s).} Data: CIFAR-10 with one fixed split; ground truth is only used to train
the teacher, for the ceiling and for testing. The weak teacher is a small CNN with 62 per cent accuracy, plus
a stronger second teacher. The student is a frozen DINO vision transformer, pretrained on ImageNet without
labels, with a small head with dropout. It is ``strong'' because with true labels it reaches 95.6 per cent. We
compare naive training, the A02 loss, a per-image alpha variant we already piloted, and our gated target, with
ablations: a random gate and different uncertainty types. Five seeds each.

\textbf{Slide 8, evaluation (70 s).} Each metric answers one question. Accuracy and PGR: does the student still
beat its teacher? The accuracy--coverage curve, our main metric: if the most uncertain images go to a human,
how accurate is the rest? AUROC of the uncertainty for errors, and separately for errors copied from the
teacher. ECE for calibration. And a diagnostic on the training set: are the images the gate refuses to harden
really the ones the teacher mislabelled? We compare with a random gate at the same q, so a gain can only come
from the uncertainty. A null result would also be informative.

\textbf{Slide 9, pilot, plan, takeaway (60 s).} We have a pilot with three seeds for the simpler per-image alpha
variant. Accuracy does not change, but error detection improves slightly and consistently, also on copied
errors, which are the hardest to detect in every row. The gated target is not run yet. First results on
December 2, ablations by December 16, poster on December 23. The main risk is that MC Dropout on a small head
gives a weak signal; the fallback is a five-head ensemble. To close: A02 makes the student confident, U09 makes
uncertainty actionable, and we test whether combining them during training gives a student that knows when
not to be trusted.

\section{Likely questions and short answers}
\Q{Why not just change $\alpha$?} Because a per-image $\alpha$ is algebraically the same as replacing $\fhat$ by
a blend that falls back to the teacher's label (Sect.\ 2.2). Changing $\fhat$ directly lets us choose a better
fallback: the student's own soft belief.

\Q{Is gating pseudo-labels new?} Not by itself. FixMatch (Sohn et al., 2020) keeps a pseudo-label only when the
confidence is high, and UPS (Rizve et al., 2021) selects pseudo-labels by uncertainty, both in semi-supervised
learning. New here: the weak-to-strong setting (every image has a weak label, some wrong, no clean labels), the
gate placed inside A02's loss, and the evaluation by abstention and copied-error detection.

\Q{Why would uncertainty know where the teacher is wrong?} It may not; that is hypothesis H3 and we measure it
directly, since we hold the true training labels for evaluation only. The intuition: teacher errors contradict
what the pretrained features suggest, so heads that fit the weak labels equally well disagree there.
Supporting evidence: in A02 (Sect.\ 5, overfitting) and in our pilot the student first learns the clean
pattern and only later fits the teacher's errors.

\Q{If the student has confidently copied an error, uncertainty is low and the gate hardens it. Then what?}
Correct, the gate cannot fix those. This is why the warm-up matters (the gate acts while disagreement still
exists) and why we report copied errors separately.

\Q{Data vs.\ model uncertainty?} Data (aleatoric): noise in the label given the image; does not vanish with more
data. Model (epistemic): disagreement between plausible models; shrinks with relevant data. Example from U09,
Table 1: passes (0.5,0.5),(0.5,0.5) are all data uncertainty; passes (1,0),(0,1) are all model uncertainty;
same total.

\Q{How strong is the student, and what was it pretrained on?} DINO ViT-S/16 (Caron et al., 2021),
self-supervised on ImageNet \emph{without labels}. Trained on true CIFAR-10 labels it reaches 95.6\,\% against
62\,\% for the teacher. Caveat: ImageNet contains similar objects, but no labels were seen; A02 uses the same
argument for its DINO experiment (no ``pretraining leakage'' of labels).

\Q{Why frozen features and a small head?} Cost: features are extracted once and each run takes seconds to
minutes, so we can afford 5 seeds and all ablations. A02's vision experiment is also a linear probe on frozen
DINO features. Limitation: results may differ with full fine-tuning.

\Q{Why CIFAR-10?} Clean ground truth is needed to count teacher errors and copied errors; CIFAR-10-C gives a
standard shift test.

\Q{What is the baseline code?} A02: \url{https://github.com/openai/weak-to-strong} (MIT licence, has a
\texttt{vision} folder for AlexNet $\to$ DINO). U09 cites no repository; MC Dropout, the decomposition and
the rejection metrics are re-implemented by us.

\Q{What are NRA, CQ, RQ?} U09's rejection metrics. Non-rejected accuracy: accuracy on what is kept.
Classification quality: share of correct decisions (kept and right, or rejected and wrong). Rejection quality:
how much misclassified cases are concentrated in the rejected set compared with the whole set.

\Q{What is PGR?} $(\text{student}-\text{weak})/(\text{ceiling}-\text{weak})$. Pilot: $(0.720-0.624)/(0.956-0.624)=0.29$.

\Q{The pilot gain is tiny (0.79 to 0.80). Is it worth it?} It is small but stable across seeds ($\pm0.004$),
with a very weak teacher and the simplest variant. It motivates the proposal; it is not the result. We will
report effect sizes with seed variation and accept a null result.

\Q{Why is ECE worse with the confidence loss in the pilot (0.041 to 0.091)?} Hardening sharpens predictions;
with a 62\,\% teacher the student then becomes over-confident. This is one more reason to gate the hardening.

\Q{Theoretical contribution?} The gradient decomposition (imitate vs.\ sharpen), the equivalence between
per-example $\alpha$ and a teacher-fallback self-target, and a one-parameter family ($q$) that contains A02 and
makes the accuracy/uncertainty trade-off measurable.

\Q{Where is the third paper (U03)?} After the 11 Nov discussion we focus on two papers, as advised; U03's idea
(uncertainty under an unseen group) survives only as the optional shift test.

\section{To check before Wednesday}
\begin{itemize}
\item The pilot table on slide 9 is copied from our shared notes (3 seeds). Whoever ran it should confirm the
numbers and that ``$\pm$'' is the standard deviation across seeds.
\item Decide who presents which slides (suggestion: one person 1--2 and 4, one 3 and 5--6, one 7--9) and rehearse
once with a timer: the script is about 1{,}250 words.
\item Submit the PDF of the slides.
\end{itemize}

\section*{References}
\small
\begin{itemize}
\item Burns, C., et al.\ (2024). Weak-to-strong generalization: Eliciting strong capabilities with weak supervision. \emph{ICML}. arXiv:2312.09390.
\item Thuy, A., \& Benoit, D.\,F.\ (2024). Explainability through uncertainty: Trustworthy decision-making with neural networks. \emph{European Journal of Operational Research}, 317, 330--340.
\item Caron, M., et al.\ (2021). Emerging properties in self-supervised vision transformers. arXiv:2104.14294.
\item Sohn, K., et al.\ (2020). FixMatch: Simplifying semi-supervised learning with consistency and confidence. \emph{NeurIPS}. arXiv:2001.07685.
\item Rizve, M.\,N., Duarte, K., Rawat, Y.\,S., \& Shah, M.\ (2021). In defense of pseudo-labeling: An uncertainty-aware pseudo-label selection framework for semi-supervised learning. \emph{ICLR}. arXiv:2101.06329.
\item Hendrycks, D., \& Dietterich, T.\ (2019). Benchmarking neural network robustness to common corruptions and perturbations. \emph{ICLR}.
\item Cited through A02, App.\ A.4 (not opened by us): Grandvalet \& Bengio (2004), Lee (2013), Reed et al.\ (2014).
\item Course slides: Hung, Y.-H.\ (2026). \emph{Trustworthy AI: Introduction to TAI} and \emph{Uncertainty Quantification} decks, NYCU.
\end{itemize}
\end{document}
